\documentclass[11pt]{article}
\usepackage[a4paper,margin=26mm,headheight=15pt]{geometry}
\usepackage[T1]{fontenc}
\usepackage{lmodern,microtype,amsmath,amssymb,amsthm,mathtools}
\usepackage{booktabs,longtable,array,enumitem,xcolor,hyperref,fancyhdr,listings}
\hypersetup{colorlinks=true,linkcolor=blue!45!black,urlcolor=blue!55!black,citecolor=blue!45!black,
pdftitle={Transcendence of e and pi in intuitionistic primitive recursive arithmetic},
pdfauthor={Research proof package},pdfsubject={Finite Hermite proofs and Lean blueprint}}
\definecolor{navy}{RGB}{26,48,76}
\definecolor{light}{RGB}{243,246,249}
\pagestyle{fancy}\fancyhf{}
\fancyhead[L]{\small Transcendence in intuitionistic PRA}
\fancyhead[R]{\small Finite Hermite proofs}
\fancyfoot[C]{\thepage}
\setlength{\parskip}{4pt}
\setlength{\parindent}{0pt}
\setlength{\emergencystretch}{3em}
\setlist{itemsep=3pt,topsep=5pt}
\lstset{basicstyle=\ttfamily\small,breaklines=true,backgroundcolor=\color{light},frame=single,rulecolor=\color{black!15},columns=fullflexible}
\newtheorem{theorem}{Theorem}[section]
\newtheorem{lemma}[theorem]{Lemma}
\newtheorem{proposition}[theorem]{Proposition}
\newtheorem{corollary}[theorem]{Corollary}
\theoremstyle{definition}
\newtheorem{definition}[theorem]{Definition}
\newtheorem{remark}[theorem]{Remark}
\newcommand{\ipra}{\mathsf{iPRA}}
\newcommand{\N}{\mathbb N}
\newcommand{\Z}{\mathbb Z}
\newcommand{\Q}{\mathbb Q}
\newcommand{\ii}{\mathrm i}
\newcommand{\tr}{\operatorname{tr}}
\newcommand{\adj}{\operatorname{adj}}
\newcommand{\Exp}{\mathcal E}
\newcommand{\norm}[1]{\left\lVert #1\right\rVert}
\newcommand{\abs}[1]{\left|#1\right|}
\newcommand{\FPint}{\mathcal I}
\newcommand{\cut}{\operatorname{cut}}
\title{\color{navy}\bfseries Transcendence of \(e\) and \(\pi\)\\
in intuitionistic primitive recursive arithmetic}
\author{A constructive proof and a Lean implementation blueprint}
\date{14 November 2026}
\begin{document}
\maketitle
\begin{abstract}
We give uniform, quantitative proofs of the transcendence of \(e\) and
\(\pi\) whose content is finite rational arithmetic with primitive recursive
bounds. For each nonzero integer polynomial \(P\), we explicitly construct a
positive integer \(B(P)\) and a precision \(N(P)\) such that the standard
primitive recursive rational names \(c_s\) of either constant satisfy
\(2B(P)|P(c_s)|>1\) whenever \(s\ge N(P)\).
This is stronger than the negative assertion that no polynomial vanishes.

The proof for \(e\) is Hermite's auxiliary-polynomial argument. The proof
for \(\pi\) replaces algebraic conjugates by a rational companion matrix and
its additive exterior powers. Newton identities replace the passage to a
splitting field. Both proofs use an explicitly chosen composite-allowed
integer in place of a prime. Polynomial integration is a finite coefficient
operation; the apparent analytic identities have exact truncated versions
and explicit remainder bounds. We give the induction audit needed to read
the argument in first-order \(\ipra\), inspect the actual ConciseIPRA
interfaces in ProveIt, and specify the remaining Lean implementation.
The accompanying Lean file is an uncompiled blueprint, not a completed
machine-checked proof.
\end{abstract}
\tableofcontents
\newpage

\section{The precise assertion and its constructive meaning}

\subsection{The theory}

We use full first-order intuitionistic logic with equality, a function
symbol for each primitive recursive function, its defining equations,
successor injectivity, \(Sx\ne0\), and induction for quantifier-free formulas.
This is the \(\ipra\) presentation implemented by
\texttt{PRAAx} in the inspected ConciseIPRA development
\cite{repoPR}. It is not Heyting arithmetic with unrestricted induction.

Integers are signed natural-number codes. Rationals are pairs consisting
of an integer numerator and a positive denominator; reduction to lowest
terms is unnecessary. Finite lists, polynomials and matrices are natural
number codes. Their equality, rational order, divisibility and all searches
used below are decidable primitive recursive operations.

\subsection{Names for the two constants}

Put
\[
 \Exp_s(z)=\sum_{j=0}^s\frac{z^j}{j!},\qquad
 e_s=\Exp_{s+1}(1),
\]
and put
\[
 A_s(x)=\sum_{j=0}^s\frac{(-1)^j x^{2j+1}}{2j+1},
 \qquad
 \pi_s=16A_s(1/5)-4A_s(1/239).
 \tag{1.1}
\]
All these are rational when their arguments are rational. The estimates
proved below give
\[
 |e-e_s|\le 2^{-s},\quad |\pi-\pi_s|\le 2^{-s},
 \qquad 2\le e_s\le3,\quad 3<\pi_s<4.
 \tag{1.2}
\]
Here and throughout, a real or complex expression denotes its explicitly
regulated rational approximations. There is no sort of arbitrary real
numbers in the object theory. More precisely, the finite Cauchy estimates
\(|e_t-e_s|,|\pi_t-\pi_s|\le2^{-s}\) hold for all \(t\ge s\).
These finite versions will suffice for the internal proof.

For \(P(X)=\sum_{j=0}^d a_jX^j\in\Z[X]\), define the positive integer
\[
 \lambda(P)=1+\sum_{j=1}^d j|a_j|4^{j-1}.
 \tag{1.3}
\]
Factoring \(x^j-y^j\) shows, by a finite sum, that
\[
 |P(x)-P(y)|\le \lambda(P)|x-y|\qquad (|x|,|y|\le4).
 \tag{1.4}
\]

\begin{theorem}[Uniform finite transcendence]\label{thm:main}
There are explicitly specified primitive recursive functions
\(B_e,B_\pi:\text{polynomial codes}\to\N_{>0}\). If \(P\ne0\), \(c=e\)
or \(c=\pi\), and
\[
 B=B_c(P),\qquad N_c(P)=4B\lambda(P)+1,
 \tag{1.5}
\]
then \(\ipra\) proves the rational inequality
\[
 \forall s\ge N_c(P)\qquad 2B\,|P(c_s)|>1.
 \tag{1.6}
\]
The theorem is uniform in the code of \(P\); it is not merely a family of
separate arguments for externally fixed polynomials.
\end{theorem}

The algorithms for \(B_e\) and \(B_\pi\) appear in
Sections~\ref{sec:e} and~\ref{sec:pi}. Their numerical size is deliberately
generous. Efficiency is not needed for primitive recursiveness.

An entirely first-order rendering of \(P(c)=0\) is
\[
 \forall s\quad |P(c_s)|\le\lambda(P)2^{-s}.
 \tag{1.7}
\]
For the named Cauchy sequence this says exactly that the polynomial value
is zero. Equation (1.6), instantiated at (1.5), refutes (1.7), since
\(2^N\ge N+1>4B\lambda(P)\).
Thus the constructive conclusion is an explicit rational separation
witness, and hence the negation of polynomial vanishing. Clearing a common
positive denominator extends the conclusion from integer to rational
polynomials.

\subsection{Removal of zero coefficients}

Trim trailing zeros by a scan through the input list. If the list is not
identically zero, find its first nonzero coefficient by another bounded
scan and write
\[
 P(X)=X^vQ(X),\qquad Q(0)\ne0.
\]
If \(Q\) is constant, set \(B_c(P)=1\): all our named values are at least
\(2\), so \(|P(c_s)|\ge1\), which directly implies (1.6).
Otherwise run the construction below on \(Q\)
and use the resulting \(B\) also for \(P\). Multiplication by \(c^v\)
or \(c_s^v\) cannot weaken a lower bound. In (1.5), use \(\lambda(P)\)
of the original input. No test of equality of arbitrary real numbers
occurs.

\section{Finite arithmetic tools}\label{sec:finite}

\subsection{The common factorial estimate}

\begin{lemma}[A prescribed coprime modulus]\label{lem:modulus}
Let \(M,A,R\ge1\) be integers. Set
\[
 L=4R^2+4AR+1,\quad q=ML,\quad p=2q+1.
 \tag{2.1}
\]
Then \(p>1\), \(p\) is coprime to \(M\), and
\[
 \frac{AR^p}{(p-1)!}<\frac14.
 \tag{2.2}
\]
There is no primality requirement on \(p\).
\end{lemma}
\begin{proof}
The equality \(p-2LM=1\) supplies an explicit inverse of \(M\) modulo
\(p\). The last \(q\) factors in \((2q)!\) are at least \(q\), so
\[
 (2q)!\ge q^q,\qquad
 \frac{AR^{2q+1}}{(2q)!}
 \le AR\left(\frac{R^2}{q}\right)^q
 \le \frac{AR}{4^q}.
\]
Induction gives \(4^q\ge q+1\). Since \(q>4AR\), the last quantity
is less than \(1/4\). Each displayed inequality is an integer inequality
after multiplication by positive denominators.
\end{proof}

A repeated arithmetic fact will be useful. For \(h\in\Z[X]\),
expansion at an integer \(j\) gives
\[
 h^{(r)}(j)=r!\,[X^r]h(j+X).
 \tag{2.3}
\]
Consequently \(h^{(r)}(j)/(p-1)!\) is an integer for \(r\ge p-1\),
and a multiple of \(p\) for \(r\ge p\).
This uses the factor \(p\) in \(r!/(p-1)!\), not the primality of \(p\).

\subsection{Complex arithmetic without square roots}

A Gaussian rational is a pair \(a+\ii b\), \(a,b\in\Q\). Use
\[
 |a+\ii b|_1=|a|+|b|.
\]
It is subadditive and submultiplicative. For a finite matrix use the maximum
row sum with these entry norms, denoted \(\|\cdot\|\).
Then
\[
 \|UV\|\le\|U\|\|V\|,\qquad |\tr U|_1\le D\|U\|
 \tag{2.4}
\]
for \(D\) rows. Maximum, row sums, and absolute values are rational
operations. On real scalars the norm is ordinary absolute value.

\subsection{Exponential bounds with a primitive recursive cutoff}

For an integer \(R\ge1\), if \(s+2\ge2R\), comparison of consecutive
terms gives, for every finite \(t>s\),
\[
 \sum_{j=s+1}^t\frac{R^j}{j!}
 \le \frac{2R^{s+1}}{(s+1)!}.
 \tag{2.5}
\]
Every subsequent ratio is at most \(1/2\), and the finite
geometric sum is at most \(2\). Also
\[
 \Exp_s(R)\le3^R. \tag{2.6}
\]
Indeed, \(j!\ge2^{j-1}\) for \(j\ge1\) gives
\(\Exp_s(1)\le3\), and the multinomial expansion of
\(\Exp_s(1)^R\) contains every term of \(\Exp_s(R)\).

A convenient, intentionally loose cutoff is
\[
 \cut(R,k)=4(R+1)^2+k+4. \tag{2.7}
\]
For \(s\ge\cut(R,k)\), both \(R^s/s!\) and the right side of (2.5)
are at most \(2^{-k}\). Let \(K=4(R+1)^2\).
Then \(K!\ge(K/2)^{K/2}\), so
\(R^K/K!\le(2R^2/K)^{K/2}\le1\); all later ratios are at most
\(1/2\). This proves the claim with spare powers of \(2\).
The same inequalities apply to matrices of norm at most \(R\).

We write \(\exp U\) for this regulated sequence of matrix polynomials.
Multiplication and addition identities are proved from finite coefficient
identities followed by (2.5). For commuting \(U,V\), the binomial
coefficient identity gives \(\exp(U+V)=\exp U\exp V\).
No completeness theorem supplies an unspecified limit.

\subsection{Integration means a rational coefficient operation}

For a polynomial \(h(t)=\sum h_jt^j\) with scalar or matrix rational
coefficients, define
\[
 \FPint(h)=\sum_j\frac{h_j}{j+1}.
 \tag{2.8}
\]
This is what \(\int_0^1h(t)\,dt\) will mean. Direct differentiation of
monomials proves
\(\FPint(h')=h(1)-h(0)\) and polynomial integration by parts.
In particular,
\[
 \FPint\bigl(t^u(1-t)^v\bigr)=\frac{u!\,v!}{(u+v+1)!}
 \quad(u,v\in\N).
 \tag{2.9}
\]
The case \(v=0\) is (2.8); integration by parts gives
\(I_{u,v}=vI_{u+1,v-1}/(u+1)\) for \(v>0\).
Thus (2.9) is an equality of rationals proved by induction.
Its value lies in \([0,1]\). Expanding in the basis
\(t^u(1-t)^v\) therefore justifies all integral estimates below by a
finite triangle inequality.

\subsection{The exact finite Hermite identity}\label{sec:finitehermite}

Let \(f\in\Q[X]\) have degree at most \(\ell\), and put
\[
 F=f+f'+\cdots+f^{(\ell)}.
\]
Then \(F-F'=f\). If \(B\) is a rational matrix, define
\[
 \begin{split}
 J_s(B)&=\FPint\!\left(B\Exp_s((1-t)B)f(tB)\right),\\
 K_s(B)&=\FPint\!\left(
       \frac{(1-t)^sB^{s+1}}{s!}F(tB)\right).
 \end{split}
\]
Polynomial differentiation gives the exact identity
\[
 F(0)\Exp_s(B)-F(B)=J_s(B)-K_s(B). \tag{2.10}
\]
In detail, the derivative of
\(\Exp_s((1-t)B)F(tB)\) is
\[
 -B\Exp_s((1-t)B)f(tB)
 +\frac{(1-t)^sB^{s+1}}{s!}F(tB);
\]
all factors commute because they are polynomials in \(B\).
Equation (2.8) proves (2.10).

Write \(|h|_b=\sum|h_j|b^j\). If \(\|B\|\le b\), (2.9) implies
\[
 \|J_s(B)\|\le b\,3^b|f|_b,\qquad
 \|K_s(B)\|\le\frac{b^{s+1}}{s!}|F|_b. \tag{2.11}
\]
More precisely, if \(F(X)=\sum F_jX^j\), then
\[
 K_s(B)=\sum_j F_j B^{s+1+j}\frac{j!}{(s+j+1)!};
\]
this also yields the stronger denominator \((s+1)!\) in (2.11).
The weaker displayed bound is sufficient.
Consequently the customary shorthand
\[
 F(0)\exp B-F(B)
 =\int_0^1B\exp((1-t)B)f(tB)\,dt
 \tag{2.12}
\]
has a specified finite meaning. It is not an appeal to a general
fundamental theorem of calculus in \(\ipra\).

\section{The constant \texorpdfstring{\(\pi\)}{pi} and Euler's identity}\label{sec:pi-name}

The alternating terms defining \(A_s(x)\) decrease for \(0\le x\le1/5\).
Grouping adjacent terms in a finite tail gives
\[
 |A(x)-A_s(x)|\le\frac{x^{2s+3}}{2s+3}.
\]
Hence
\[
 |\pi-\pi_s|
 \le\frac{20}{(2s+3)5^{2s+3}}\le2^{-s}. \tag{3.1}
\]
The same bound holds between stages \(t\ge s\).
The elementary inequalities \(x-x^3/3\le A_s(x)\le x\)
give \(3<\pi_s<4\) by rational arithmetic. Similarly
\(|e-e_s|\le2/(s+2)!\le2^{-s}\).

We give the algebra behind Euler's identity.
For the formal coefficient list
\(A(X)=X-X^3/3+X^5/5-\cdots\), put
\[
 R(X)=\frac{1+\ii X}{1-\ii X}.
\]
The identity
\[
 \exp(2\ii A(X))=R(X) \tag{3.2}
\]
is a coefficient identity. Both sides have constant coefficient \(1\) and
satisfy \((1+X^2)Y'=2\ii Y\).
If \(Y=\sum y_jX^j\), its coefficients are uniquely determined by
\[
 y_1=2\ii,\qquad
 (j+1)y_{j+1}=2\ii y_j-(j-1)y_{j-1}\quad(j\ge1).
 \tag{3.3}
\]
Thus (3.2) to any specified degree follows by finite rational induction.

Here is an explicit justification for evaluation. At \(\rho=1/2\),
\[
 \sum_{j=0}^s\frac{\rho^{2j+1}}{2j+1}\le\frac23.
\]
The absolute coefficient sum of
\(\Exp_L(2\ii A_s(X))\), weighted by \(\rho^j\), is at most
\(\Exp_L(4/3)\le9\). If \(J\le\min(L,2s+1)\), its coefficients through
degree \(J\) equal those of \(R(X)\). For \(0\le x\le1/5\),
\[
 \left|\Exp_L(2\ii A_s(x))-R(x)\right|_1
 \le 12(2/5)^{J+1}. \tag{3.4}
\]
The polynomial tail is at most \(9(2x)^{J+1}\); the geometric
tail of \(R(X)=1+2\sum_{j\ge1}(\ii X)^j\) is at most
\(2x^{J+1}/(1-x)\). This proves (3.4) by finite majorants.
The same reasoning applies with \(-\ii\).

The exact Gaussian rational identity
\[
 (12+5\ii)^4=\ii(28560+239\ii),\qquad 13^4=28561
 \tag{3.5}
\]
implies
\[
 R(1/5)^4R(1/239)^{-1}=\ii.
\]
Square and use exponential addition:
\[
 \exp(\ii\pi)=R(1/5)^8R(1/239)^{-2}=-1.
 \tag{3.6}
\]
In completely finite form, for every \(k\),
\[
 t\ge k+30,\quad h\ge\cut(4,k+30)
 \quad\Longrightarrow\quad
 |\Exp_h(\ii\pi_t)+1|_1\le2^{-k}. \tag{3.7}
\]
Here is a completely finite error calculation.
Take \(m=\max(h,2t+2)\), \(z_1=2\ii A_t(1/5)\), and
\(z_2=-2\ii A_t(1/239)\).
Apply (3.4) with \(L=m\), \(J=2t+1\), to the eight factors
\(\Exp_m(z_1)\) and the two factors \(\Exp_m(z_2)\).
Each factor and its rational target have norm at most \(3\).
A telescoping product bounds the error by
\(120\,3^9(2/5)^{2t+2}\).
The finite multinomial expansion of
\(\Exp_m(z_1)^8\Exp_m(z_2)^2\) agrees through total degree \(m\)
with \(\Exp_m(8z_1+2z_2)=\Exp_m(\ii\pi_t)\).
Since \(8|z_1|_1+2|z_2|_1\le16/5+4/239<4\), its remaining
terms have absolute sum at most \(2^{-k-30}\), by (2.5)--(2.7).
The difference between \(\Exp_m(\ii\pi_t)\) and
\(\Exp_h(\ii\pi_t)\) is at most another \(2^{-k-30}\).
Thus the total bound is
\[
 120\,3^9(2/5)^{2t+2}+2\cdot2^{-k-30}<2^{-k}.
\]
This proves (3.7) with finite operations, uniformly in \(t,k\).

For identification with the conventional circle constant, (1.1) is
Machin's formula. Its elementary verification uses
\(\tan(2\arctan(1/5))=5/12\),
\(\tan(4\arctan(1/5))=120/119\), and
\(\tan(4\arctan(1/5)-\arctan(1/239))=1\).
The resulting angle lies in the principal arctangent interval, so it
equals \(\arctan1=\pi/4\).
This identifies the name with ordinary \(\pi\); internally only the
rational definition, (3.1) and (3.7) are used.



\section{Transcendence of \texorpdfstring{\(e\)}{e}}\label{sec:e}

After the reduction in Section 1, write
\[
 Q(X)=\sum_{j=0}^n a_jX^j,\qquad n\ge1,\quad a_0a_n\ne0,
 \qquad S=\sum_{j=0}^n|a_j|.
\]
Use the following positive integers:
\[
 \begin{gathered}
 R=(2n)^{n+1},\quad A=n3^{2n}S,\quad M=|a_0|n!,\\
 L=4R^2+4AR+1,\quad q=ML,\quad p=2q+1,\\
 \ell=(n+1)p-1,\qquad
 D=\ell!\bigl((n+1)!\bigr)^p.
 \end{gathered} \tag{4.1}
\]
Define
\[
 h(X)=X^{p-1}\prod_{j=1}^n(X-j)^p,\qquad
 f(X)=\frac{h(X)}{(p-1)!},\qquad F=\sum_{r=0}^{\ell}f^{(r)}.
 \tag{4.2}
\]

\subsection{A nonzero integer}

At \(j=1,\ldots,n\), derivatives of order \(r<p\) vanish, and all
higher derivatives are divisible by \(p\) after division by \((p-1)!\).
At zero, derivatives below \(p-1\) vanish and
\[
 f^{(p-1)}(0)=(-1)^{np}(n!)^p.
\]
Thus all \(F(j)\) are integers and
\[
 F(j)\equiv0\pmod p\ (1\le j\le n),\qquad
 F(0)\equiv(-1)^{np}(n!)^p\pmod p. \tag{4.3}
\]
Therefore
\[
 T=\sum_{j=0}^na_jF(j)
 \equiv a_0(-1)^{np}(n!)^p\pmod p
\]
is a nonzero integer. To justify this even for composite \(p\), the
product on the right divides \(M^p\) up to sign, while
\(M\) has the explicit inverse \(-2L\) modulo \(p\).
Hence it is invertible modulo \(p\). In particular \(|T|\ge1\).

\subsection{A small error}

The scalar instance of (2.12), rescaled to \([0,j]\), is
\[
 e^jF(0)-F(j)=J(j),\qquad
 J(j)=\int_0^j e^{j-t}f(t)\,dt. \tag{4.4}
\]
For a coefficient bound, evaluate absolute coefficients of \(h\) at \(n\):
\[
 |h|_n=n^{p-1}\prod_{j=1}^n(n+j)^p\le R^p.
\]
Expanding in \(t^u(j-t)^v\), the rational beta-integral estimate gives
\[
 |J(j)|\le n3^n\frac{R^p}{(p-1)!}.
\]
Consequently Lemma~\ref{lem:modulus}, with the slightly larger \(A\)
chosen in (4.1), gives
\[
 \left|\sum_{j=0}^na_jJ(j)\right|<\frac14. \tag{4.5}
\]
Multiplying (4.4) by \(a_j\) and summing yields
\[
 F(0)Q(e)=T+\sum_j a_jJ(j),\qquad
 |F(0)Q(e)|>\frac34. \tag{4.6}
\]
The sum of the absolute coefficients of \(h\) is
\(\bigl((n+1)!\bigr)^p\). Each contribution to \(F(0)\) is
a coefficient multiplied by \(r!/(p-1)!\), whence
\[
 |F(0)|\le D,\qquad
 |Q(e)|>\frac{3}{4D}>\frac1{2D}. \tag{4.7}
\]
Set
\[
 \boxed{B_e(P)=2D}. \tag{4.8}
\]

\subsection{The complete finite version}

Write \(u_s=\Exp_s(1)\) and define the rational coefficient integrals
\[
 \begin{split}
 J_s(j)&=\int_0^j \Exp_s(j-t)f(t)\,dt,\\
 K_s(j)&=\int_0^j (j-t)^sF(t)/s!\,dt.
 \end{split}
\]
Polynomial differentiation gives
\[
 F(0)\Exp_s(j)-F(j)=J_s(j)-K_s(j). \tag{4.9}
\]
The multinomial expansion gives, for \(j\le n\), \(s+2\ge2n\),
\[
 0\le u_s^j-\Exp_s(j)\le\frac{2n^{s+1}}{(s+1)!}.
 \tag{4.10}
\]
The terms of total degree at most \(s\) agree; the remaining positive
terms are bounded by the finite exponential tail at \(j\).

Put
\[
 H_F=(\ell+1)Dn^\ell,\quad W=Sn(H_F+2D),\quad R_1=2n,\quad
 t_0=4R_1^2+4W+1,\quad s_0=2t_0. \tag{4.11}
\]
Each coefficient of \(F\) has absolute value at most \(D\):
for \(h(X)=\sum h_aX^a\), a coefficient of \(f^{(r)}\)
has multiplier \(a!/((a-r)!(p-1)!)\le\ell!\), and the contributions
to a fixed coefficient sum over distinct \(a\).
Consequently \(|F|_n\le H_F\).
The coefficient triangle inequality or the beta-integral formula gives
\[
 \left|\sum a_jJ_s(j)\right|<\frac14,
\]
and
\[
 \left|\sum a_jK_s(j)\right|
 +\left|F(0)\sum a_j(u_s^j-\Exp_s(j))\right|
 \le W\frac{R_1^s}{s!}<\frac14
 \qquad(s\ge s_0). \tag{4.12}
\]
For the last strict inequality, at \(s=2t_0\) use
\((2t_0)!\ge t_0^{t_0}\), giving
\(W(R_1^2/t_0)^{t_0}\le W4^{-t_0}<1/4\).
The ratio of successive terms is at most \(1/2\) thereafter.
Equations (4.9)--(4.12) put \(F(0)Q(u_s)\) within \(1/2\)
of \(T\). Thus
\[
 s\ge s_0\ \Longrightarrow\ |Q(u_s)|>\frac1{2D}.
 \tag{4.13}
\]
Since our name is \(e_t=u_{t+1}\), the safe threshold \(t\ge s_0\)
suffices. This establishes the separation using only finite rational
arithmetic; Section~\ref{sec:arithmetization} transfers it to the common
precision (1.5).

\section{Finite linear algebra for the \texorpdfstring{\(\pi\)}{pi} proof}\label{sec:linear}

\subsection{Companion matrices and approximate eigenvectors}

For \(q(X)=cX^n+b_{n-1}X^{n-1}+\cdots+b_0\), \(c>0\), take
\[
 A=\begin{pmatrix}
 0&1&0&\cdots&0\\
 0&0&1&\cdots&0\\
 \vdots&&&\ddots&\vdots\\
 0&0&0&\cdots&1\\
 -b_0/c&-b_1/c&-b_2/c&\cdots&-b_{n-1}/c
 \end{pmatrix}.
 \tag{5.1}
\]
For \(v(z)=(1,z,\ldots,z^{n-1})^T\),
\[
 (A-zI)v(z)=-\frac{q(z)}c\,e_n. \tag{5.2}
\]
This exact identity holds for every Gaussian rational \(z\); it does
not assume that a root exists.

\subsection{Additive exterior powers}

Index the basis of the \(k\)-th exterior power by increasing \(k\)-element
subsets of \(\{1,\ldots,n\}\). Define \(A^{[k]}\) by
\[
 A^{[k]}(v_1\wedge\cdots\wedge v_k)
  =\sum_{j=1}^k v_1\wedge\cdots\wedge Av_j
                       \wedge\cdots\wedge v_k.
 \tag{5.3}
\]
For \(k=0\), this is the \(1\)-by-\(1\) zero matrix.
Entries are integer linear combinations of entries of \(A\);
wedge signs are finite permutation parities. If \(\|A\|\le L\),
\[
 \|A^{[k]}\|\le kL. \tag{5.4}
\]
For a fixed row subset, its off-diagonal entries replace one selected
index, and its diagonal is the sum of the selected diagonal entries.
The row sum is bounded by the sum of the corresponding \(k\) row sums
of \(A\). Define
\[
 B=\bigoplus_{k=0}^n A^{[k]},\qquad D_0=2^n. \tag{5.5}
\]
In particular \(B\) has a specified zero block.

\begin{lemma}[Exterior exponential identity]\label{lem:exterior}
For the regulated matrix exponential,
\[
 \tr(\exp B)=\det(I+\exp A). \tag{5.6}
\]
\end{lemma}
\begin{proof}
Finite multinomial induction on \(r\) gives
\[
 (A^{[k]})^r(v_1\wedge\cdots\wedge v_k)
 =\sum_{r_1+\cdots+r_k=r}
 \frac{r!}{r_1!\cdots r_k!}
 A^{r_1}v_1\wedge\cdots\wedge A^{r_k}v_k.
\]
Divide by \(r!\) and compare coefficients to obtain
\(\exp(A^{[k]})=\bigwedge^k(\exp A)\).
For any matrix \(U\), the trace of \(\bigwedge^k U\) is the sum of its
principal \(k\)-minors. Expanding the determinant by selecting an identity
entry or a \(U\) entry in each row gives
\(\det(I+U)=\sum_{k=0}^n\tr(\bigwedge^k U)\).
Combine these finite coefficient identities and the tail bound (2.5).
\end{proof}

For later use, an especially simple fully finite error bound is
\[
 \left|\det(I+\Exp_s(A))-\tr(\Exp_s(B))\right|
 \le C_{\rm ext}2^{-s-1},\qquad
 C_{\rm ext}=n!\bigl(1+3^{2L}\bigr)^n. \tag{5.7}
\]
Indeed the two polynomials in an extra parameter \(z\),
\(\det(I+\Exp_s(zA))\) and \(\tr(\Exp_s(zB))\), agree through degree
\(s\). The second has no terms of higher degree.
The absolute coefficient majorant of the first, evaluated at \(z=2\),
is at most \(C_{\rm ext}\): use the determinant's \(n!\) signed
products and \(\Exp_s(2L)\le3^{2L}\).
Every term of degree exceeding \(s\) is smaller at \(z=1\) by at
least \(2^{-s-1}\). This proves (5.7) directly.

\subsection{A trace identity replacing algebraic conjugates}

\begin{lemma}[Removing the zero factor]\label{lem:newton}
Let \(M\) be a \(D_0\)-by-\(D_0\) integer matrix. Write, by scanning
its characteristic coefficients,
\[
 \chi_M(X)=X^mG(X),\qquad G\in\Z[X]\ \hbox{monic},\quad G(0)\ne0.
\]
Let \(C_G\) be the integer companion matrix of \(G\), including the empty
matrix if \(G=1\). Every rational polynomial \(u\) satisfies
\[
 \tr u(M)=m\,u(0)+\tr u(C_G). \tag{5.8}
\]
\end{lemma}
\begin{proof}
For \(r\ge1\), Newton's coefficient recurrence gives
\(\tr M^r=\tr C_G^r\); the traces at degree zero differ by \(m\).
Here is a derivation of that recurrence needing no roots.
The polynomials \(d_M(t)=\det(I-tM)\) and
\(d_C(t)=\det(I-tC_G)\) are equal because \(\chi_M=X^mG\).
Differentiate the Leibniz determinant formula and use the adjugate
identity and
\[
 (I-tM)\sum_{j=0}^r t^jM^j=I-t^{r+1}M^{r+1}.
\]
To every prescribed degree this gives
\[
 -t\,d_M'(t)=d_M(t)\sum_{j\ge1}\tr(M^j)t^j.
\]
The constant coefficient of \(d_M\) is \(1\); coefficient comparison
recursively determines each trace without division by an unknown
quantity. The same recurrence holds for \(C_G\). Linear combination
proves (5.8). All calculations are truncated polynomial identities.
\end{proof}

The companion identity \(G(C_G)=0\) follows by multiplying out the
companion matrix, or by coefficient comparison in the adjugate identity.
Thus this section uses neither a splitting field nor any theorem
asserting the existence of complex roots.

\section{The matrix Hermite lemma}\label{sec:matrix-hermite}

\begin{theorem}[Effective nonvanishing of an exponential trace]
\label{thm:matrix}
Let \(B\) be a rational \(D_0\)-by-\(D_0\) matrix, \(D_0\ge1\).
Suppose \(c\ge1\) is an integer, \(M=cB\) has integer entries, and
\[
 \chi_M=X^mG,\qquad m\ge1,\quad G(0)\ne0.
\]
Let \(r=\deg G\), \(H_G=\sum_{j=0}^r|g_j|\), and choose an integer
\(b\ge\max(1,\|B\|)\). Define
\[
 \begin{gathered}
 V=cb,\quad R=H_GV^{r+1},\quad A_0=D_0b3^b,\quad
 M_0=mc|G(0)|,\\
 L_0=4R^2+4A_0R+1,\quad q_0=M_0L_0,\quad p=2q_0+1,\\
 \ell=(r+1)p-1,\qquad U=\ell!\,c^\ell H_G^p.
 \end{gathered} \tag{6.1}
\]
Then
\[
 \left|\tr(\exp B)\right|>\frac{3}{4U}. \tag{6.2}
\]
\end{theorem}
\begin{proof}
Put
\[
 h(Y)=Y^{p-1}G(Y)^p,\qquad
 f(X)=\frac{h(cX)}{(p-1)!},\qquad F=\sum_{j=0}^{\ell}f^{(j)}.
 \tag{6.3}
\]
For \(j<p\), \(h^{(j)}\) is divisible by \(G\): differentiating
\(G^p\) fewer than \(p\) times leaves at least one of its \(p\) factors
undifferentiated in each product-rule term.
Thus \(h^{(j)}(C_G)=0\).

For \(j\ge p\), all coefficients of
\(h^{(j)}/(p-1)!\) are integer multiples of \(p\).
At a coefficient of degree \(a-j\), the multiplier is
\[
 \frac{a!}{(a-j)!(p-1)!}
 =\binom aj\frac{j!}{(p-1)!},
\]
which is integral and divisible by \(p\). Since
\[
 f^{(j)}(C_G/c)=\frac{c^j h^{(j)}(C_G)}{(p-1)!},
\]
\(\tr F(C_G/c)\) is an integer divisible by \(p\).
At zero the only possible nonzero derivative below \(p\) is
\(f^{(p-1)}(0)=c^{p-1}G(0)^p\); all later contributions
are integer multiples of \(p\).
Apply (5.8) to \(u(Y)=F(Y/c)\). It follows that
\[
 T=\tr F(B)\in\Z,\qquad
 T\equiv m c^{p-1}G(0)^p\pmod p. \tag{6.4}
\]
The right side divides \(M_0^p\) up to sign.
Since \(p-2L_0M_0=1\), it is invertible modulo \(p\).
Hence \(T\ne0\) and \(|T|\ge1\).
This argument includes \(G=1\), using the empty companion matrix.

For the estimate,
\[
 |f|_b
 =\frac{|h|_{cb}}{(p-1)!}
 \le\frac{V^{p-1}(H_GV^r)^p}{(p-1)!}
 \le\frac{R^p}{(p-1)!}.
\]
Take the trace in (2.12). By (2.4), (2.11) and
Lemma~\ref{lem:modulus},
\[
 F(0)\tr(\exp B)=T+J,\qquad
 |J|\le\frac{A_0R^p}{(p-1)!}<\frac14. \tag{6.5}
\]
Thus \(|F(0)\tr(\exp B)|>3/4\).
The absolute coefficient sum of \(h(cX)\) is at most \(c^\ell H_G^p\),
so \(|F(0)|\le\ell!c^\ell H_G^p=U\).
This proves (6.2).
\end{proof}

\begin{remark}
The integer \(m\) is a multiplicity read from an integer polynomial.
It is not a number of roots obtained by deciding equality of complex
numbers. No eigenvalue or algebraic conjugate is selected.
\end{remark}

\section{Transcendence of \texorpdfstring{\(\pi\)}{pi}}\label{sec:pi}

Let the stripped input be
\[
 P_0(X)=\sum_{j=0}^d a_jX^j,\qquad d\ge1,\quad a_0a_d\ne0.
\]
Define
\[
 q(Z)=P_0(-\ii Z)P_0(\ii Z). \tag{7.1}
\]
This belongs to \(\Z[Z]\): odd coefficients cancel; each even
coefficient is a sum of integer products with signs.
Its degree is \(n=2d\), its leading coefficient is \(c=a_d^2>0\), and
\[
 q(\ii x)=P_0(x)P_0(-x). \tag{7.2}
\]
No vanishing equation is assumed.

Let \(A\) be (5.1) for \(q\). Use the explicit integer bound
\[
 L=\max\left(1,\sum_{j=0}^{n-1}|[Z^j]q|\right). \tag{7.3}
\]
Then \(\|A\|\le L\). Form \(B\) by (5.5), let \(D_0=2^n\), and
take \(b=nL\). The matrix \(cB\) is integral and \(\|B\|\le b\).
Its specified zero block makes the coefficient scan in
\(\chi_{cB}=X^mG\) return \(m\ge1\).
Compute \(U\) from (6.1). Equations (5.6) and (6.2) give
\[
 |\det(I+\exp A)|>\frac{3}{4U}. \tag{7.4}
\]

\subsection{Relating the determinant to the polynomial value}

Put \(\alpha=\ii\pi\), \(v=(1,\alpha,\ldots,\alpha^{n-1})^T\),
and \(W=I+\exp A\). The difference-of-powers identity yields
\[
 (\exp A-\exp\alpha\,I)v
 =\int_0^1\exp((1-t)A)\exp(t\alpha)(A-\alpha I)v\,dt.
 \tag{7.5}
\]
One may verify this by comparing coefficients in
\(A^j-\alpha^jI=(A-\alpha I)\sum_{u=0}^{j-1}A^u\alpha^{j-1-u}\);
the beta-integral formula gives the coefficients \(1/j!\).
The finite exponential calculus therefore justifies (7.5).

Using (3.6), (5.2), \(|\alpha|_1\le4\), and \(c\ge1\), we obtain
\[
 \|Wv\|\le3^{L+4}|q(\alpha)|_1. \tag{7.6}
\]
Also \(\|W\|\le T_0=1+3^L\).
Each cofactor of \(W\) is a signed sum of \((n-1)!\) products,
so a row sum of \(\adj W\) is at most
\[
 K_{\rm adj}=n!T_0^{n-1}.
\]
The first coordinate of \(v\) is \(1\). Taking that coordinate in
\(\adj W\,Wv=(\det W)v\) gives
\[
 |\det W|\le K|q(\alpha)|_1,\qquad
 K=K_{\rm adj}3^{L+4}. \tag{7.7}
\]
No inverse of \(W\) is assumed.

Finally let \(H=\sum_{j=0}^d|a_j|4^j\).
Equation (7.2) gives
\(|q(\alpha)|_1=|P_0(\pi)P_0(-\pi)|\le H|P_0(\pi)|\).
Combining with (7.4),
\[
 |P_0(\pi)|>\frac{3}{4UKH}>\frac1{2UKH}.
 \tag{7.8}
\]
For the original input put
\[
 \boxed{B_\pi(P)=2UKH}. \tag{7.9}
\]
These are explicitly given primitive recursive integers.

\subsection{A direct finite separation threshold}\label{sec:pi-finite}

We now supply all cutoffs needed to prove the rational version directly.
In addition to the constants above, put
\[
 \begin{aligned}
 V_F&=(\ell+1)U b^\ell,&
 k_H&=16D_0bV_F+1,\\
 k_E&=16U K_{\rm adj}4^{n-1}+1,&
 t_\pi&=k_E+30,\\
 h&=\max\{\cut(b,k_H),\,16UC_{\rm ext}+1,\,
                       \cut(4,k_E+30)\}.
 \end{aligned} \tag{7.10}
\]
Each coefficient of \(F\) is at most \(U\) in absolute value, by the
same factorial-coefficient argument used for (4.11), so
\(|F|_b\le V_F\).
Equations (2.10)--(2.11) give
\[
 |F(0)\tr\Exp_h(B)-T|
 \le|\tr J_h|+|\tr K_h|
 <\frac14+\frac1{16}=\frac5{16}.
\]
Indeed \(|\tr K_h|\le D_0bV_F\,2^{-k_H}<1/16\).
As \(|T|\ge1\) and \(|F(0)|\le U\),
\[
 |\tr\Exp_h(B)|>\frac{11}{16U}.
\]
By (5.7) and \(h\ge16UC_{\rm ext}+1\),
\[
 |\det(I+\Exp_h(A))|>\frac{10}{16U}. \tag{7.11}
\]

Take any \(t\ge t_\pi\), and set \(z=\ii\pi_t\).
The exact finite identity
\[
 (\Exp_h(A)-\Exp_h(z)I)v(z)
 =\left(\sum_{j=1}^h\frac1{j!}
                    \sum_{u=0}^{j-1}A^uz^{j-1-u}\right)
                    (A-zI)v(z)
 \tag{7.12}
\]
has norm at most \(3^{L+4}|q(z)|_1\).
For completeness, each term with \(u+v=j-1\) has
\(1/j!\le1/(u!v!)\); summing is bounded by
\(\Exp_h(L)\Exp_h(4)\le3^{L+4}\).
The extra term \((1+\Exp_h(z))v(z)\) has norm at most
\(4^{n-1}2^{-k_E}\), by (3.7).
The adjugate estimate therefore gives
\[
 |\det(I+\Exp_h(A))|
 \le K H|P_0(\pi_t)|+K_{\rm adj}4^{n-1}2^{-k_E}.
\]
Its last term is less than \(1/(16U)\).
Subtract it from (7.11):
\[
 KH|P_0(\pi_t)|>\frac9{16U}>\frac1{2U},
 \qquad
 t\ge t_\pi\Longrightarrow
 |P_0(\pi_t)|>\frac1{B_\pi(P)}.
 \tag{7.13}
\]
This is a proof involving only rational matrices, polynomials and
finite coefficient sums. It closes the analytic shorthand in (7.8)
without any appeal to an unspecified limit.



\section{Reading the argument inside \texorpdfstring{\(\ipra\)}{iPRA}}
\label{sec:arithmetization}

\subsection{A purely finite proof of the main theorem}

The direct finite statements (4.13) and (7.13) have already proved
\[
 t\ge t_c(P)\quad\Longrightarrow\quad
 |P(c_t)|>1/B_c(P), \tag{8.1}
\]
where \(t_e=s_0\) and \(t_\pi=k_E+30\).
The reduction \(P=X^vP_0\) preserves the inequality because
\(c_t\ge2\). Constant monomials were handled directly in Section 1.

Fix \(s\ge N=4B\lambda(P)+1\), and choose the explicit integer
\[
 t=\max(s,t_c(P)).
\]
Finite Cauchy comparison of the names and the polynomial Lipschitz
identity give
\[
 |P(c_s)|\ge |P(c_t)|-\lambda(P)|c_t-c_s|
 >\frac1B-\lambda(P)2^{-s}>\frac{3}{4B}>\frac1{2B}. \tag{8.2}
\]
For the penultimate step use \(2^s\ge s+1>4B\lambda(P)\).
This proves Theorem~\ref{thm:main} directly, without a real-number
quantifier, a completeness principle, or a convergence search.

If \(c_s=u/v\), with \(v>0\), its conclusion is the integer inequality
\[
 2B\left|\sum_{j=0}^d a_j u^jv^{d-j}\right|>v^d.
 \tag{8.3}
\]
For the pi name, \(u,v\) are computed by the two finite rational sums
in (1.1). Internally no quotient type of rationals is necessary.

\subsection{Why the bounds are primitive recursive}

All programs are fixed and uniform in the input polynomial code.
They use addition, multiplication, integer comparison, bounded scans,
factorial, exponentiation, finite list operations, coefficient convolution,
and matrix operations of a prescribed dimension.
The characteristic polynomial may be computed directly by the determinant
formula, enumerating permutations with a bound such as \(D_0^{D_0}\).
The exterior basis consists of masks below \(2^n\). The zero
multiplicity is a scan of at most \(D_0+1\) integer coefficients.
Every dimension, loop count, derivative order and Taylor length is a
displayed primitive recursive function of the input.

A symbolic expression for a very large factorial is only a readable
description of the corresponding ordinary integer. It does not change
the mathematical definition of \(B\) or \(N\).
The proof makes no claim that these deliberately loose bounds are
practical for numerical computation.

\subsection{The induction audit}

Primitive recursiveness of an algorithm alone does not prove its
correctness in \(\ipra\). Here the actual correctness arguments have
quantifier-free invariants, summarized below. Bounded families of
equalities or inequalities are represented by one primitive recursive
Boolean check. Induction on that check is quantifier-free induction;
entries, dimensions and finite bounds are free numerical parameters.

\small
\begin{longtable}{@{}p{.28\textwidth}p{.28\textwidth}p{.37\textwidth}@{}}
\toprule
Construction & Finite bound & Quantifier-free invariant\\\midrule
\endhead
List/rational arithmetic &
Input lengths and explicit codes &
Cross-multiplied numerator/denominator equations; positive denominators\\
Polynomial products and derivatives &
Coefficient degree and derivative order &
Coefficient arrays equal convolution or factorial formulas\\
Factorial domination &
The chosen \(q\), then a free finite tail endpoint &
Integer products and cross-multiplied inequalities in (2.2), (2.5)\\
Exponential products &
Taylor order and bounded compositions of that order &
Coefficientwise binomial/multinomial equality\\
Polynomial integration &
Degree and exponents \(u,v\) &
Rational beta-integral recurrence (2.9)\\
Machin identity &
Taylor order and \(\cut(R,k)\) &
Coefficient recurrence (3.3) and positive majorant bound\\
Exterior matrices &
Subsets below \(2^n\); tuples below \((r+1)^k\) &
Multinomial expansion of the \(r\)-th power\\
Determinants and adjugates &
Permutations of the fixed finite dimension &
Partial signed permutation sums and cancellation identities\\
Newton recurrence &
Characteristic degree and requested power &
Bounded prefix of coefficient equations of the formal logarithmic derivative\\
Zero multiplicity &
At most \(D_0+1\) coefficients &
Earlier coefficients vanish; the first retained coefficient does not\\
Hermite congruences &
Degree \(\ell\) and matrix dimension &
Each derivative coefficient has its asserted factorial divisor\\
Final separation &
The explicit \(h,t_c,N\) &
Cross-multiplied accumulated error inequality\\
\bottomrule
\end{longtable}
\normalsize

The Newton argument does not require strong induction: retain the
bounded prefix of earlier trace equalities in one array check.
The same device handles the coefficient recurrences and partial
sums. Nested loops do not introduce induction on an unbounded
existential or universal formula. One proves the arrays' defining
recurrences and the stated Boolean invariants. Universally quantified
theorem parameters are introduced after those derivations.

In particular there is no induction on a formula asserting that all
polynomials possess witnesses, and no use of
\(\neg\neg\exists n\,R(n)\Rightarrow\exists n\,R(n)\).
All witnesses are terms given in advance. Branches concern only
decidable integer and finite-array tests. The proof therefore uses
intuitionistic reasoning and quantifier-free arithmetic induction
throughout.

One must also avoid a different hidden strengthening: a universal
evaluator for arbitrary primitive recursive program codes is not
itself a primitive recursive function. Nothing here assumes one.
The required arithmetic functions are a fixed finite collection of
programs compiled from composition and primitive recursion.

\section{Lean formalization: a concrete completion path}\label{sec:lean}

\subsection{The repository interfaces actually inspected}

The repository snapshot is
\[
 \texttt{0bf379a84212abd05e9a3f0b1320a3d7da509cf9}.
\]
Its toolchain file specifies Lean \texttt{v4.32.0}.
We inspected the ConciseIPRA README and four source files
\cite{repoPR,repoDeduction,repoNative}:
\begin{itemize}[nosep]
\item \texttt{PR.lean} and \texttt{Deduction.lean};
\item \texttt{PRAToScalar.lean} and \texttt{NativeBiInterpretation.lean}.
\end{itemize}
The shared namespace is \texttt{LeanProofs.ConciseIPRA}.
The relevant interfaces are:

\begin{center}\small
\begin{tabular}{@{}ll@{}}
\toprule
API & Role\\\midrule
\texttt{PR k} & Syntax of primitive recursive symbols of arity \(k\)\\
\texttt{PTm n} & Terms over \(n\) free variables\\
\texttt{PRAAx} & Defining equations and quantifier-free induction\\
\texttt{Prf PRAAx} & Intuitionistic object-theory derivations\\
\texttt{PRA.qfInd} & Induction requiring an explicit \texttt{QF} proof\\
\texttt{NFm.ofPRA} & Translation to native ConciseIPRA formulas\\
\texttt{pra\_iff\_native} & Faithful transfer of iPRA derivations\\
\bottomrule
\end{tabular}
\end{center}

The directory \texttt{docs/drafts} was also inventoried. It contains
related analysis packages, including an FTA draft, but the present proof
does not depend on their claims: the matrix argument removes any need
for an effective FTA. No repository change or remote commit is part of
this delivery.

\subsection{The actual formal target}

The accompanying \texttt{lean/Blueprint.lean} defines finite rational
names, a separation predicate, and interfaces for the missing programs
and derivations. At the object level the goal is a closed statement
of the form
\begin{lstlisting}
Prf PRAAx [] (AllAccepted concretePrograms.piTest)
\end{lstlisting}
and the analogous statement for \(e\).
The test returns \(1\) on the zero polynomial, and otherwise checks
the rational separation at the explicit precision index.
The stronger all-later-precisions theorem (1.6) can be added with a
second numerical parameter; one precision already proves transcendence.

An ambient Lean theorem
\(\texttt{Transcendental}\ \Q\ \texttt{Real.pi}\), however correctly
proved, would not establish that \(\ipra\) proves the assertion.
Avoiding a \texttt{Classical} command would not suffice either.
The object derivations must use the specified quantifier-free
induction rule. Evaluating a Python program on examples also supplies
no universal derivation.

\subsection{Module order and mathematical obligations}

\begin{enumerate}
\item \textbf{Codec and elementary arithmetic.}
Build fixed PR codes for signed integers, positive-denominator rationals,
lists, polynomial arrays and matrices. Prove defining equations,
cross-multiplication, bounded conjunction and comparison lemmas.

\item \textbf{Finite polynomial calculus.}
Implement convolution, derivative, formal integration, coefficient
norms and factorial bounds. Prove (2.3), (2.9), (2.10) internally.
This supplies the main common infrastructure.

\item \textbf{Finite exponential and Machin calculations.}
Prove (2.5)--(2.7), the coefficient product identities, the arctangent
recurrence (3.3), and the rational Euler bound (3.7).
The internal proof needs no general real or complex integration library.

\item \textbf{Companion and exterior matrices.}
Use bit masks for subsets and inversion parity for wedge signs.
Prove (5.2), the multinomial formula in
Lemma~\ref{lem:exterior}, principal-minor expansion, and (5.7).

\item \textbf{Trace algebra.}
Prove the adjugate identity by the finite determinant formula, the
Newton recurrence, (5.8), and the companion annihilation identity.
This avoids formalizing an algebraic closure.

\item \textbf{The two Hermite arguments.}
Prove the derivative congruences and the explicit composite-parameter
lemma. Instantiate them for \(e\) and the integer matrix \(cB\).
Carry the magnitude bounds through the exact finite formulas,
including (4.13) and (7.13).

\item \textbf{Final arithmetic theorems.}
Define the programs computing \(B_c(P)\), \(\lambda(P)\), \(N_c(P)\)
and the rational test (8.3). Build internal proofs with
\texttt{Prf} constructors and \texttt{PRA.qfInd}. Prove the encoding
specifications as well as the external mathematical statements.

\item \textbf{Transport and verification.}
Apply \texttt{pra\_iff\_native.mp} to the completed iPRA derivations.
Print the axioms of the final Lean declarations. Serialization into a
repository Python checker format, if wanted, is a separately verified
last step; the existing interpretation theorem does not provide that
serializer automatically.
\end{enumerate}

The blueprint supplies no \texttt{axiom} or \texttt{sorry}-based
transcendence theorem. Its \texttt{Programs} and \texttt{Derivations}
records are explicit obligations to instantiate.
No Lean executable was available in this environment, so the blueprint
has not been compiled. The human mathematical proof and the exact Python
checks have different, expressly stated statuses.

\section{Artifacts, validation and scope}

The archive contains this source and compiled PDF, the Lean blueprint and
a detailed roadmap, exact-arithmetic Python checks, an executable
description of the bound construction, and a verification log.

The \(e\) checks exercise the exact polynomial identity, integer
congruences, a composite modulus \(p=25\), and finite rational separation
examples. The matrix checks exercise exterior-power coefficient
identities, the polynomial transformation, companion residuals,
Newton traces, and Hermite congruences for rational and nilpotent matrices,
including composite moduli \(9\) and \(25\).
Additional finite-analysis checks verify the Gaussian Machin identities,
Cauchy estimates on representative indices, and exact scalar and matrix
Hermite truncation defects. These are implementation checks of finite
identities, not universal iPRA proof certificates.

The bound-construction artifact keeps enormous final factorials and
powers in a symbolic arithmetic expression tree. This is a presentation
choice: the formulas define ordinary primitive recursive integers.
Example outputs do not claim that those astronomical integers or their
corresponding Taylor expansions were materialized.

The mathematical machinery is the Hermite method, organized to make
its finite arithmetic content explicit. Existing Coq and Isabelle
formalizations provide useful comparisons and implementation experience
\cite{bernard,eberl}. Their completed proofs are not imported as proofs
of internal \(\ipra\)-derivability. No historical novelty claim is made
for the matrix organization.

\begin{thebibliography}{9}
\bibitem{repoPR}
Vladimir Reshetnikov, \emph{ProveIt: ConciseIPRA, PR.lean}.
Snapshot \texttt{0bf379a84212abd05e9a3f0b1320a3d7da509cf9}.
\href{https://github.com/VladimirReshetnikov/ProveIt/blob/0bf379a84212abd05e9a3f0b1320a3d7da509cf9/Logic/PrimitiveRecursiveArithmetic/ConciseIPRA/Lean/ConciseIPRA/PR.lean}{Pinned source on GitHub}.

\bibitem{repoDeduction}
Vladimir Reshetnikov, \emph{ConciseIPRA, Deduction.lean}, same snapshot.
\href{https://github.com/VladimirReshetnikov/ProveIt/blob/0bf379a84212abd05e9a3f0b1320a3d7da509cf9/Logic/PrimitiveRecursiveArithmetic/ConciseIPRA/Lean/ConciseIPRA/Deduction.lean}{Pinned source on GitHub}.

\bibitem{repoNative}
Vladimir Reshetnikov, \emph{ConciseIPRA, NativeBiInterpretation.lean},
same snapshot.
\href{https://github.com/VladimirReshetnikov/ProveIt/blob/0bf379a84212abd05e9a3f0b1320a3d7da509cf9/Logic/PrimitiveRecursiveArithmetic/ConciseIPRA/Lean/ConciseIPRA/NativeBiInterpretation.lean}{Pinned source on GitHub}.

\bibitem{bernard}
S. Bernard, Y. Bertot, L. Rideau and P.-Y. Strub,
\emph{Formal Proofs of Transcendence for \(e\) and \(\pi\) as an
Application of Multivariate and Symmetric Polynomials}.
Certified Programs and Proofs, 2016; preprint, 2015.
\url{https://arxiv.org/abs/1512.02791}.

\bibitem{eberl}
M. Eberl, \emph{The Hermite--Lindemann--Weierstrass Transcendence
Theorem}. Archive of Formal Proofs, 2021.
\url{https://isa-afp.org/entries/Hermite_Lindemann.html}.
\end{thebibliography}
\end{document}

